% SPDX-License-Identifier: Apache-2.0
\section{Two Clocks, One Run}
\label{sec:x-twoclock}

A unit's process waits on one clock, so a design of two clocks is a
unit of units. \code{Two} holds a counter on the default clock and a
counter on \code{Slow}, which takes six ticks a cycle where the
default takes two, and names no clock itself. Each child names its
own, in its register and at both its ports, and the clock a port is
on is the second argument of its type: \code{Out<U<8>, Slow>}.

Nothing crosses between the two domains here, and that is the point of
the example rather than an omission. What it is for is the checking.
A crossing would put the part worth watching inside a \code{Crossing},
where a testbench cannot see it; two counters put it at the ports,
where it can.

\lstinputlisting[rangebeginprefix=//\ begin\{,rangebeginsuffix=\},%
rangeendprefix=//\ end\{,rangeendsuffix=\},includerangemarker=false,%
linerange=clock-clock,caption={\code{lib/examples/ex\_twoclock.rs}:
the second clock, which is three ticks of the first to one of
its.},label={lst:x-slowclock}]{lib/examples/ex_twoclock.rs}

Until this the build could not check such a unit at all, which is
issue 131, and the reason is worth stating because it is the kind of
thing a test suite hides rather than shows. The generated testbench
advanced in cycles of the first clock and compared every output one of
those cycles after driving it. For a port on a clock three times
slower that is a value its own edge has not produced yet, so the two
agreed for part of every period and disagreed for the rest.

The fix is that a port carries its clock. The clock is written in the
type and nowhere else by the time the netlist exists, so \code{\#[lower]}
takes it from there into the lowering, the ports file names it after
the port's width, and the generator reads it and checks each port one
cycle of its own clock after driving it.

That this works is not something the run can show by itself, since a
testbench that is wrong in the same direction as the design agrees
with it. What shows it is breaking it on purpose: with the offset put
back to the first clock's period, the slow port differs at 5000,
11000, 17000 and every six thousand after, which is that clock's
period, and the co-simulation fails. With it right, both simulators
agree with the run for the whole of it.

The slow counter's input tells a second story of the same kind, which
is issue 405. The testbench applies every input before every tick of
its own clock, and a port on the slower clock samples only the last
application before its own edge. The generator read that value from
the trace at the tick plus the port's period, which for the default
clock is the next edge and for \code{Slow} is a tick up to four
after it: an input the run changed between two of the slow edges was
offered to the entity before the edge it belonged after, and the
entity counted an edge early. Nothing in the tree changed an input
between a slow clock's edges until a line buffer on a pixel clock was
handed its column by the run, so this too was invisible. The
generator now reads the value at the port's next edge, which is the
same tick as before for the default clock and the right one for every
other. The run here is the check: \code{go\_slow} is told on and off
by turns one tick after each of the slow edges, so what the counter
must count is the value in force at the edge and not the one after.
With the old rule the slow count differs at 6950, 8950 and 10950,
and at the same three offsets every twelve thousand after, which is
every other period of its clock; with the new one both simulators
agree with the run.

\lstinputlisting[rangebeginprefix=//\ begin\{,rangebeginsuffix=\},%
rangeendprefix=//\ end\{,rangeendsuffix=\},includerangemarker=false,%
linerange=units-units,caption={\code{lib/examples/ex\_twoclock.rs}:
the two counters and the parent that holds
them.},label={lst:x-twoclock}]{lib/examples/ex_twoclock.rs}
